\begin{remark}\label{fga-rem-4-3}\dcref{4.3}\source{fga-explained:page-0079}\uses{fga-def-4-2}
For a covering family $\{U_i\to U\}$, the category of descent data is independent,
up to canonical equivalence, of the chosen fiber products $U_{ij}$ and $U_{ijk}$.
With a cleavage, an object $\xi\in\mathcal X(U)$ gives the pullbacks $\xi_i$ and
the canonical transition isomorphisms on $U_{ij}$; the cocycle condition follows
from the uniqueness of cartesian arrows.  The cleavage-free definition uses
cartesian cubes over the diagram of the $U_i$, $U_{ij}$, and $U_{ijk}$, and is
canonically equivalent to the usual descent category.
\end{remark}

\begin{remark}\label{fga-rem-4-4}\dcref{4.4}\source{fga-explained:page-0081}\uses{fga-def-4-2}
The descent construction can be expressed using the covering sieve $h_U$.  If
$F:h_U\to\mathbf{Set}$ is a morphism of fibered categories, its values on
$U_i$, $U_{ij}$, and $U_{ijk}$ form descent data.  This gives a functor
\[
 \operatorname{Hom}(h_U,\mathcal X)\longrightarrow
 \mathcal X_{\mathrm{desc}}(\{U_i\to U\}),
\]
independent of a cleavage; after choosing one, it agrees with pullback from
$\mathcal X(U)$ through the $2$-Yoneda equivalence.
\end{remark}

\begin{remark}\label{fga-rem-4-10}\dcref{4.10}\source{fga-explained:page-0084}\uses{fga-def-4-8}
We use ``prestack'' for a fibered category whose presheaves of morphisms are
sheaves, and ``stack'' when, in addition, every descent datum for objects is
effective.  This is the terminology in the source, although ``separated
prestack'' is sometimes used for the first condition.
\end{remark}

\begin{proposition}\label{fga-pro-4-20}\dcref{4.20}\source{fga-explained:page-0089}\uses{fga-def-3-31,fga-def-4-8}
Let $\mathcal X\to\mathcal C$ be fibered and let $\mathcal X^{\mathrm{cart}}$
be its associated fibered groupoid.  If $\mathcal X$ is a stack, then
$\mathcal X^{\mathrm{cart}}$ is a stack.  If $\mathcal X$ is a prestack and
$\mathcal X^{\mathrm{cart}}$ is a stack, then $\mathcal X$ is a stack.
\end{proposition}
\begin{proof}
The two fibered categories have the same objects and their cartesian arrows are
the isomorphisms.  Thus descent data have the same objects and effective objects;
it remains to prove the prestack assertion.  For objects $\xi,\eta\in
\mathcal X(U)$, a compatible family of isomorphisms after a covering glues
uniquely because $\mathcal X$ is a prestack.  If an arrow $a:\xi\to\eta$ becomes
cartesian after pullback, it is an isomorphism locally; a local inverse glues to
an inverse of $a$.  Hence $a$ is cartesian, proving the claim.
\end{proof}

\begin{lemma}\label{fga-lem-4-22}\dcref{4.22}\source{fga-explained:page-0091}\uses{fga-lem-2-61}
For a ring map $A\to B$ and an $A$-module $M$, the sequence
\[
0\longrightarrow M\xrightarrow{a_M}B\otimes_A M
 \mathrel{\substack{\longrightarrow\\[-.6ex]\longrightarrow}}
B\otimes_A B\otimes_A M\tag{4.22.1}
\]
is exact, where $a_M(m)=1\otimes m$ and the two arrows are
$b\otimes m\mapsto b\otimes1\otimes m$ and
$b\otimes m\mapsto1\otimes b\otimes m$.
\end{lemma}
\begin{proof}
The assertion is local on $A$.  After tensoring the equalizer diagram with the
faithfully flat $B$, the equalizer is the standard equalizer of the two maps
$B\otimes_A B\rightrightarrows B\otimes_A B\otimes_A B$, with splitting induced
by multiplication in the first two factors.  Faithful flatness reflects exactness,
so the original equalizer is $M$ via $a_M$.
\end{proof}

\begin{remark}\label{fga-rem-4-24}\dcref{4.24}\source{fga-explained:page-0092}\uses{fga-rem-2-56}
For the wild flat topology, quasi-coherent sheaves need not even form a
prestack.  The cover by local schemes at all closed points and the direct sum of
their structure sheaves give compatible projections on every member of the cover,
but no global morphism inducing all those projections.  Thus morphisms fail the
sheaf condition.
\end{remark}

\begin{lemma}\label{fga-lem-4-25}\dcref{4.25}\source{fga-explained:page-0093}\uses{fga-prop-4-12,fga-lem-4-17}
Let $S$ be a scheme and $\mathcal X$ a fibered category over $(\mathrm{Sch}/S)$.
Assume that $\mathcal X$ is a Zariski stack and that, for every faithfully flat
map of affine $S$-schemes $V\to U$, the functor
$\mathcal X(U)\to\mathcal X(V\to U)$ is an equivalence.  Then $\mathcal X$
is an fpqc stack.
\end{lemma}
\begin{proof}
The prestack condition is local, so it follows from the Zariski condition and
the affine hypothesis.  For effectiveness, first replace a cover by a single
faithfully flat arrow using disjoint unions and the refinement lemma.  If the
target is affine, cover the source by quasi-compact opens that remain surjective;
the affine hypothesis and the prestack condition glue the resulting objects.
For a general target, restrict to an affine open cover, glue there by the Zariski
stack property, and use uniqueness on overlaps to satisfy the cocycle condition.
The glued object pulls back to the prescribed descent datum, proving effectivity.
\end{proof}

\begin{lemma}\label{fga-lem-4-26}\dcref{4.26}\source{fga-explained:page-0093}\uses{fga-def-4-8}
The category $\mathcal X(\varnothing)$ is equivalent to the terminal category:
between any two of its objects there is a unique arrow.
\end{lemma}
\begin{proof}
The empty scheme has the empty Zariski covering.  Descent for this covering has
one object and one arrow, and the Zariski stack condition identifies it with
$\mathcal X(\varnothing)$.
\end{proof}

\begin{lemma}\label{fga-lem-4-27}\dcref{4.27}\source{fga-explained:page-0093}\uses{fga-def-4-8}
If $U=\coprod_iU_i$ is a disjoint union of open subschemes, restriction induces
an equivalence
\[
\mathcal X(U)\xrightarrow{\sim}\prod_i\mathcal X(U_i).
\]
\end{lemma}
\begin{proof}
Full faithfulness is the sheaf condition for the presheaf of arrows.  Given
objects $\xi_i\in\mathcal X(U_i)$, the intersections $U_i\cap U_j$ are empty
for $i\ne j$; use Lemma~\ref{fga-lem-4-26} to supply the unique transition data
and apply Zariski descent.  This gives an inverse up to canonical isomorphism.
\end{proof}

\begin{theorem}\label{fga-the-4-29}\dcref{4.29}\source{fga-explained:page-0097}\uses{fga-thm-4-23}
The fibered category $(\operatorname{QCohComm}/S)$ of quasi-coherent
$\mathcal O$-algebras over $(\operatorname{Sch}/S)$ is a stack for the fpqc
topology.
\end{theorem}
\begin{proof}
On an affine fpqc cover $\operatorname{Spec}B\to\operatorname{Spec}A$, a
descent datum on a quasi-coherent algebra is a $B$-algebra with a compatible
isomorphism after tensoring with $B\otimes_AB$.  The equalizer module of
Lemma~\ref{fga-lem-4-22} descends the underlying module; the unit and
multiplication descend because their diagrams commute after the faithfully flat
base change.  The resulting module is therefore an $A$-algebra and recovers the
given datum.  The quasi-coherent sheaf stack theorem and locality of algebra
maps glue these affine constructions.
\end{proof}

\begin{proposition}\label{fga-pro-4-35}\dcref{4.35}\source{fga-explained:page-0102}\uses{}
Consider a commutative diagram of schemes
\[
\begin{array}{ccccc}
X&\xrightarrow{f}&Y&\xrightarrow{g}&Z\\[-1mm]
\downarrow\!\xi&&\downarrow\!\eta&&\downarrow\!\zeta\\[-1mm]
U&\xrightarrow{\phi}&V&\xrightarrow{\psi}&W .
\end{array}
\]
For an $\mathcal O_Z$-module $\mathcal L$, the canonical base-change and
pullback maps satisfy
\[
\xi_*\alpha_{f,g}(\mathcal L)\,\beta_{f,\phi}(g^*\mathcal L)
 \beta_{g,\zeta}(\mathcal L)
=\beta_{gf,\xi}(\mathcal L)\,\alpha_{\phi,\psi}(\zeta_*\mathcal L).
\]
Equivalently, the natural square of $\mathcal O_U$-modules formed by these maps
commutes.
\end{proposition}
\begin{proof}
Evaluate both composites on a section $s$ of $\zeta_*\mathcal L$ over an open
subset of $W$.  Each sends the pullback of $s$ to the same section of
$(gf)^*\mathcal L$ over the inverse image in $X$; equality of sheaf maps follows.
\end{proof}

\begin{proposition}\label{fga-pro-4-37}\dcref{4.37}\source{fga-explained:page-0102}\uses{fga-pro-4-35}
In a cartesian square
\[
\begin{array}{ccc}
X&\xrightarrow{f}&Y\\[-1mm]
\downarrow\!\xi&&\downarrow\!\eta\\[-1mm]
U&\xrightarrow{\phi}&V .
\end{array}
\]
assume that $\eta$ is proper and of finite presentation, and that $\mathcal L$
is quasi-coherent, of finite presentation and flat over $V$.  If
$H^1(Y_v,\mathcal L_v)=0$ for every $v\in V$, then $\eta_*\mathcal L$ is locally
free and the base-change map
\[
\beta_{f,\phi}(\mathcal L):\phi^*\eta_*\mathcal L
\xrightarrow{\sim}\xi_*f^*\mathcal L
\]
is an isomorphism.
\end{proposition}
\begin{proof}
The assertion is local on $V$.  For noetherian $V$, it is cohomology and base
change: properness and flat finite presentation, together with the vanishing of
$H^1$, give local freeness of $\eta_*\mathcal L$ and the stated isomorphism.
For arbitrary $V$, write an affine neighbourhood as a filtered limit of
noetherian affine schemes and spread out $(Y,\mathcal L)$; semicontinuity gives
an open neighbourhood on which the first cohomology vanishes.  The noetherian
result descends to that neighbourhood, and these neighbourhoods cover $V$.
\end{proof}

\begin{proposition}\label{fga-pro-4-41}\dcref{4.41}\source{fga-explained:page-0108}\uses{fga-pro-4-35}
For a cartesian square
\[
\begin{array}{ccc}
X&\xrightarrow{f}&Y\\[-1mm]
\downarrow\!\xi&&\downarrow\!\eta\\[-1mm]
U&\xrightarrow{\phi}&V .
\end{array}
\]
let $\mathcal L_\eta$ be an $\mathcal O_Y$-module,
$\mathcal M_\eta=\eta_*\mathcal L_\eta$, and
$\mathcal L_\xi=f^*\mathcal L_\eta$, $\mathcal M_\xi=\xi_*\mathcal L_\xi$.
Then the two canonical composites
\[
(\phi\xi)^*\mathcal M_\eta\longrightarrow\mathcal L_\xi,
\qquad
(\phi\xi)^*\mathcal M_\eta\longrightarrow\xi^*\mathcal M_\xi
 \xrightarrow{\tau_\xi}\mathcal L_\xi,
\]
where the first uses $f^*\tau_\eta$ and the second uses base change, coincide.
\end{proposition}
\begin{proof}
For a section $s$ of $\mathcal L_\eta$ over the inverse image of an open subset
of $V$, both composites send $(\phi\xi)^*s$ to $f^*s$.  Hence they agree on all
sections and therefore as morphisms of sheaves.
\end{proof}

\begin{lemma}\label{fga-lem-4-47}\dcref{4.47}\source{fga-explained:page-0113}\uses{fga-def-4-42,fga-lem-4-44}
If $X\to Y$ is a $G$-torsor, then
\[
\delta'_a:G\times G\times X\longrightarrow X\times_YX\times_YX,
\qquad(g,h,x)\longmapsto(ghx,hx,x),
\]
is an isomorphism.
\end{lemma}
\begin{proof}
After pulling back along a cover on which the torsor is trivial, the map is the
identity under the standard identifications
$X\times_YX\simeq G\times X$ and
$X\times_YX\times_YX\simeq G\times G\times X$.  Isomorphy is fpqc local, so
the original map is an isomorphism.
\end{proof}

\begin{remark}\label{fga-rem-4-48}\dcref{4.48}\source{fga-explained:page-0114}\uses{fga-thm-4-46}
An algebraic space over $S$ is an étale sheaf on $(\operatorname{Sch}/S)$ that
is étale-locally represented by a scheme.  Algebraic spaces and algebraic stacks
provide effective fppf descent in situations where descent for schemes fails;
properties local for the étale topology (such as flatness, smoothness, and the
Cohen--Macaulay condition) extend directly, while global properties such as
properness require a separate descent argument.
\end{remark}
